\section{Conclusion}
\label{sec:conclusion}

We presented an empirical comparison of six parameter-efficient continual fine-tuning strategies
for mitigating catastrophic forgetting in small language models (1.5B--3.8B parameters) under a
strict compute budget of a single NVIDIA T4 GPU.
Using QLoRA as a shared foundation, we evaluated Full Fine-Tuning, Vanilla LoRA, LoRA+EWC,
LoRA+Replay, O-LoRA, and our proposed Minimal Mixed Rehearsal (MMR) on a four-task GLUE sequence
(SST-2\,$\to$\,MRPC\,$\to$\,CoLA\,$\to$\,MNLI) across two instruction-tuned architectures
(Qwen2.5-1.5B-Instruct and Phi-3.5-mini-instruct).
MMR achieves \todo{the best average accuracy--compute trade-off}, reducing forgetting substantially
relative to Vanilla LoRA while incurring \todo{X\%} of the wall-clock overhead of O-LoRA and
requiring no Fisher information computation or frozen parameter snapshots.
Our results further show that orthogonality-based methods are most effective when replay data
cannot be retained, while EWC provides limited additional benefit over Vanilla LoRA at the
four-task scale studied here.
We release all code, prompt templates, and result JSON files to support full reproducibility.

Future work should pursue three natural extensions.
First, scaling MMR to larger models (7B+) would test whether the method's simplicity is preserved
when adapter rank must increase to match model capacity, and whether its Pareto advantage over
O-LoRA widens or narrows with scale.
Second, extending the evaluation to generative tasks---abstractive summarization, open-ended
question answering, and code generation---would test MMR's robustness beyond the classification
prompt paradigm; generative forgetting may require new metrics beyond sequence-level accuracy,
such as ROUGE overlap on held-out summaries or execution-based pass rates.
Third, dynamic replay-ratio scheduling---adapting $f$ at each task boundary based on an online
estimate of the current forgetting rate---could further reduce the number of replay examples
required while maintaining retention, moving toward a truly data-minimal, parameter-free continual
learning recipe suited for privacy-sensitive deployment.
